\documentclass[11pt,a4paper]{article}

\usepackage{amsmath,amssymb,amsthm}
\usepackage[utf8]{inputenc}
\usepackage[T1]{fontenc}
% \usepackage{lmodern}  % not installed in sandbox; default fonts used
\usepackage[hidelinks]{hyperref}
\usepackage{booktabs}
\usepackage{enumitem}
\usepackage[margin=2.5cm]{geometry}
\usepackage{listings}

\lstset{
  basicstyle=\ttfamily\small,
  breaklines=true,
  columns=fullflexible,
  showstringspaces=false,
  keepspaces=true,
  aboveskip=1ex,
  belowskip=1ex,
}

\newtheorem{theorem}{Theorem}[section]
\newtheorem{proposition}[theorem]{Proposition}
\newtheorem{corollary}[theorem]{Corollary}
\newtheorem{remark}[theorem]{Remark}

\newcommand{\AWs}{\mathcal{A}_W^\star}
\newcommand{\AWP}{\mathcal{A}_W^{\mathrm{PNT}}}
\newcommand{\CW}{\mathcal{C}_W}
\newcommand{\DW}{\mathcal{D}_W}
\newcommand{\SW}{\mathcal{S}_W}
\newcommand{\Theta}{\Theta_W}
\newcommand{\What}{\widehat{W}}
\newcommand{\bphi}{\varphi}

\title{TS6 --- Lean 4 Formalization Layer\\
for the TS Programme on the Smoothed Prime Variance}
\author{Serge Durand}
\date{\today}

\begin{document}
\maketitle

\begin{abstract}
This short note is the dedicated formalization layer of the TS programme
(TS1--TS5) on the empirically-centered variance of smoothed prime sums in
arithmetic progressions. TS6 formalizes, in Lean 4 over Mathlib4, the four
purely algebraic identities that constitute the exact kernel of the
programme: the abstract Parseval identity on a finite abelian group, its
specialization to $(\mathbb Z/q\mathbb Z)^\times$, the centering-shift
identity, and the Rankin--Selberg decomposition. The effective layer
(TS3, TS4) is scaffolded conditionally on two named hypothesis structures
(large sieve, effective Chebyshev) that Mathlib4 does not yet provide; no
\texttt{axiom} is used anywhere. TS6 is delivered here in YELLOW status:
every Lean statement is precisely typed and its proof plan is documented,
but the proof-term bodies remain \texttt{sorry} because the sandbox in
which this note was produced does not permit installation of a Lean
toolchain. Closing the exact-kernel \texttt{sorry}s to reach GREEN is a
routine exercise of a few hours in a live Mathlib4 environment. This note
states the scope, the architecture, the target theorems, the current
limits, and the reproducibility instructions.
\end{abstract}

\section{Scope}\label{sec:scope}

The TS programme (TS1--TS5) organizes results about the quantity
\[
\AWs(x,q)
 \;:=\; \sum_{\substack{a = 1 \\ (a,q) = 1}}^{q}
  \bigl|\theta_W(x; q, a) - \bar\theta_W(x, q)\bigr|^{2},
\]
the empirically-centered variance of the smoothed prime sum
$\theta_W(x; q, a) = \sum_{p \equiv a \,(q)} (\log p)\, W(p/x)$ across
invertible residue classes modulo $q$. The programme's three layers are:

\begin{itemize}[leftmargin=2em,nosep]
\item[\textbf{Exact.}] $\AWs(x,q) = \CW(x,q)$ for every integer
  $q \ge 2$, and the centering-shift identity relating $\AWP$ to $\CW$.
\item[\textbf{Effective.}] $\sum_{q \le Q,\, q\text{ prime}} \AWs(x,q)
  \le 2\lVert W\rVert_\infty^{2}\,(Q^{2}+2x)\, x\log(2x)$ via the large
  sieve and Chebyshev.
\item[\textbf{Open.}] Three asymptotic questions, precisely scoped.
\end{itemize}

TS6 targets \emph{only} the exact layer of that classification. The aim is
to obtain a Lean 4 formal certificate that the four exact identities hold,
as a final piece of documentation for the programme.

\smallskip

\noindent\textbf{What this note does not do.}
TS6 is deliberately narrow. It does not formalize the empirical
observations of TS1 \S5 or TS5 \S3 (regression fits, the factor-5
diagonal-overshoot observation). It does not attempt to close the effective
layer: the large sieve inequality and the effective Chebyshev bound
$\vartheta(y) \le 2y$ are not presently in Mathlib4 in the required form,
and TS6 refuses to discharge them by \texttt{axiom}. It does not formalize
any of the three open problems of TS5 \S5. These exclusions are stated
here once and enforced throughout.

\section{Exact kernel in Lean 4}\label{sec:exact}

The central targets are named TS6-A through TS6-D.

\subsection*{TS6-A --- Abstract Parseval on a finite abelian group}

Let $G$ be a finite abelian group and $f : G \to \mathbb C$. Define the
empirical mean $\mu(f) = |G|^{-1} \sum_{g} f(g)$, the variance
$\mathrm{Var}(f) = \sum_{g} |f(g) - \mu(f)|^{2}$, the Fourier coefficient
$\widehat f(\chi) = \sum_{g} f(g)\,\overline{\chi(g)}$, and the
non-principal character-side mean square
$\widehat C(f) = |G|^{-1} \sum_{\chi \ne 1} |\widehat f(\chi)|^{2}$.

\begin{theorem}[TS6-A, abstract Parseval identity]\label{thm:A}
For every finite abelian group $G$ and every $f : G \to \mathbb C$,
\[
\mathrm{Var}(f) \;=\; \widehat C(f).
\]
\end{theorem}

The proof is the standard one-line expansion using Schur orthogonality on
$G$. The Lean target is the file \texttt{TS6/Exact/FiniteCharacterVariance.lean},
which imports \texttt{Mathlib.NumberTheory.MulChar.Basic} and uses the
lemma \texttt{MulChar.sum\_apply\_eq\_zero\_iff} for non-trivial characters.

\subsection*{TS6-B --- Modular specialization}

Specializing Theorem~\ref{thm:A} to $G = (\mathbb Z/q\mathbb Z)^\times$
and $f(r) = \theta_W(x; q, r)$ (viewed as a function of a class) yields:

\begin{theorem}[TS6-B, exact identity for general $q$, TS2 Thm 2.1]
\label{thm:B}
For every integer $q \ge 2$, every $x > 0$, and every bounded weight $W$
with compact support in $(0, \infty)$,
\[
\AWs(x, q) \;=\; \CW(x, q).
\]
\end{theorem}

The Lean target is \texttt{TS6/Exact/DirichletVariance.lean}. The proof
is a direct application of TS6-A through the bijection
$\mathrm{MulChar}\,(\mathbb Z/q)^\times\,\mathbb C \;\longleftrightarrow\;
\mathrm{DirichletCharacter}\,\mathbb C\, q$.

\subsection*{TS6-C --- Centering-shift identity}

The centering-shift identity compares the empirical-mean centering to the
PNT-predicted centering $x \What(1)/\bphi(q)$.

\begin{theorem}[TS6-C, centering-shift identity, TS2 Prop 2.2]
\label{thm:C}
Under the hypotheses of Theorem~\ref{thm:B},
\[
\AWP(x, q) - \CW(x, q)
  \;=\; \frac{\bigl|\Theta(x, \chi_0) - x\,\What(1)\bigr|^{2}}{\bphi(q)}.
\]
\end{theorem}

The proof is pure shift-of-origin for the variance of a finite sequence:
$\sum_{r} |t_r - \mu|^{2} = \sum_{r} |t_r - \bar t|^{2} + |G|\,|\bar t - \mu|^{2}$,
applied at $\mu = x \What(1)/\bphi(q)$ and combined with Theorem~\ref{thm:B}.
The Lean target is \texttt{TS6/Exact/CenteringShift.lean}.

\subsection*{TS6-D --- Rankin--Selberg decomposition}

Writing $a_p = (\log p)\,W(p/x)$ and
$T_{0}(x, q) = \sum_{(p, q) = 1} a_{p}$, define the diagonal
$\DW(x, q) = \sum_{(p, q) = 1} |a_{p}|^{2}$ and the off-diagonal
\[
\SW_{q}(x) \;=\; \sum_{\substack{p_{1} \ne p_{2} \\
  p_{1} \equiv p_{2}\,(q) \\ (p_{1} p_{2}, q) = 1}}
  a_{p_{1}}\,\overline{a_{p_{2}}}.
\]

\begin{theorem}[TS6-D, Rankin--Selberg decomposition, TS5 Prop 3.7]
\label{thm:D}
For every integer $q \ge 2$,
\[
\AWs(x, q) \;=\; \DW(x, q) - \frac{|T_{0}(x, q)|^{2}}{\bphi(q)}
  + \SW_{q}(x).
\]
\end{theorem}

This identity is the algebraic output of expanding $\CW$ via the
character-sum form of Theorem~\ref{thm:B} and splitting the resulting
prime-pair double sum into $p_{1} = p_{2}$ (diagonal) vs.
$p_{1} \ne p_{2}$ (off-diagonal) contributions, subtracting the
principal-character square. The Lean target is
\texttt{TS6/Structure/RankinSelbergDecomp.lean}. No analytic input is
used: the identity is structural.

\section{Modular specialization in code}\label{sec:modular}

The core of TS6-B in Lean has the following shape (simplified):

\begin{lstlisting}
variable (q : N) [NeZero q] (a : N -> C) (S : Finset N)

noncomputable def theta (r : (ZMod q)^x) : C :=
  sum n in S, if (n : ZMod q) = (r : ZMod q) then a n else 0

noncomputable def theta_bar : C :=
  (sum r : (ZMod q)^x, theta a q S r) / (totient q : C)

noncomputable def AWstar : R :=
  sum r : (ZMod q)^x, norm (theta a q S r - theta_bar a q S)^2

noncomputable def CW : R :=
  (1 / (totient q : R)) *
    sum chi in (univ : Finset (DirichletCharacter C q)).filter (. ne 1),
      norm (sum n in S, a n * chi n)^2

theorem TS6_B_exact_identity
  (hq : 2 <= q) : AWstar a q S = CW a q S := by
  -- specialize parseval_identity to G = (ZMod q)^x
  -- carry coefficients through the MulChar ~ DirichletCharacter bijection
  sorry
\end{lstlisting}

The single \texttt{sorry} is the specialization step: every piece is in
place to push Theorem~\ref{thm:A} through the isomorphism and rename.

\section{Centering-shift in code}\label{sec:shift}

TS6-C is packaged as a separate module because it depends on TS6-B but is
otherwise pure linear algebra over $\mathbb C$. The core shift-of-origin
lemma is this:

\begin{lstlisting}
-- For any finite-indexed complex sequence (t_i) and any target mu in C,
--   sum i, ||t i - mu||^2
--     = sum i, ||t i - (sum j, t j)/|I|||^2  +  |I| * ||mu - (avg)||^2.
lemma variance_shift_of_origin
  {I : Type*} [Fintype I] (t : I -> C) (mu : C) :
  sum i, norm (t i - mu)^2 =
    sum i, norm (t i - (sum j, t j) / (Fintype.card I : C))^2
    + (Fintype.card I : R) *
        norm (mu - (sum j, t j) / (Fintype.card I : C))^2 := by
  ring_nf  -- after expanding norms via star multiplication
  sorry
\end{lstlisting}

TS6-C follows by instantiating this at
$I = (\mathbb Z/q)^\times$, $t_r = \theta_W(x; q, r)$, and
$\mu = x \What(1)/\bphi(q)$, then using TS6-B to replace the first sum on
the right by $\CW(x, q)$ and computing the second term as
$|\Theta(x, \chi_0) - x \What(1)|^{2}/\bphi(q)$.

\section{Current limits of formalization}\label{sec:limits}

TS6 does \emph{not} claim to have closed any Lean proof in the current
session: \textbf{the Lean 4 toolchain is not available in the environment
where this paper is produced}. Every statement in the Lean files
\texttt{TS6/Exact/FiniteCharacterVariance.lean},
\texttt{TS6/Exact/DirichletVariance.lean},
\texttt{TS6/Exact/CenteringShift.lean}, and
\texttt{TS6/Structure/RankinSelbergDecomp.lean} is typed and its proof
sketch is documented in a comment block, but the proof-term bodies are
\texttt{sorry}. No \texttt{axiom} is used anywhere. The statements are
proof-plan-complete; the tactic scripts to close them are routine in a
live Mathlib4 environment and are marked as future work.

The effective-layer statements in
\texttt{TS6/Effective/AveragedBounds.lean} are conditional on two typed
hypothesis structures \texttt{LargeSieveInequality} and
\texttt{ChebyshevThetaBound} --- not \texttt{axiom}s --- and will remain
so until those inputs land in Mathlib4.

\smallskip

\noindent\textbf{Summary table.}
\begin{center}
\begin{tabular}{lllc}
\toprule
Target & TS1--5 source & Lean file & Current status \\
\midrule
TS6-A & kernel for TS1, TS2 & \texttt{Exact/FiniteCharacterVariance} & YELLOW \\
TS6-B & TS2 Thm 2.1 & \texttt{Exact/DirichletVariance} & YELLOW \\
TS6-C & TS2 Prop 2.2 & \texttt{Exact/CenteringShift} & YELLOW \\
TS6-D & TS5 Prop 3.7 & \texttt{Structure/RankinSelbergDecomp} & YELLOW \\
TS6-E & TS3 Thm 1.1 & \texttt{Effective/AveragedBounds} & YELLOW (by design) \\
TS6-F & TS4 Thm 1.1 & \texttt{Effective/AveragedBounds} & YELLOW (by design) \\
TS6-G & TS4 Cor 1.2 & \texttt{Effective/AveragedBounds} & YELLOW (by design) \\
\bottomrule
\end{tabular}
\end{center}

YELLOW means: statement typed, proof plan written, no \texttt{axiom},
\texttt{sorry} explicit. The target for a subsequent session with a Lean
toolchain available is GREEN on TS6-A through TS6-D. TS6-E through
TS6-G stay YELLOW until Mathlib supplies the two missing analytic
inputs; they will not become GREEN in any near-term TS6 iteration.

\section{Reproducibility / build instructions}\label{sec:repro}

To verify or extend this formalization layer in a live Lean environment:

\begin{enumerate}[nosep,leftmargin=2em]
\item Install \texttt{elan} and Lean 4 per the
  \href{https://leanprover-community.github.io/}{leanprover-community}
  install guide.
\item Create a Mathlib-based Lean project with \texttt{lake new TS6 math}.
\item Copy the \texttt{lean/TS6/} directory of this bundle into
  \texttt{TS6/TS6/} inside the new project.
\item Add \texttt{import TS6.Exact.FiniteCharacterVariance} (and the three
  other main files) to the project's root file.
\item Run \texttt{lake build}. Each \texttt{sorry} will be flagged; closing
  them is the formalization work.
\end{enumerate}

The hash manifest \texttt{TS6\_FROZEN.txt} in the bundle records the
SHA-256 hash of each file in TS6 as released.

\section*{Labels and honesty}

All claims in this note are labelled by their logical type, consistent with
the TS programme's discipline:

\begin{itemize}[nosep,leftmargin=2em]
\item Theorems TS6-A, TS6-B, TS6-C, TS6-D are \textbf{exact identities}.
  Their proofs in the paper setting are part of TS1--TS5 and are
  independently verified there at machine precision on 672 prime-$q$ and
  126 composite-$q$ triples (see TS1 \S5 and TS2 \S4 of the merged
  bundle). Their Lean formalization status is YELLOW as described.
\item TS6-E, TS6-F, TS6-G are \textbf{effective theorems} whose
  paper-side proofs are in TS3 and TS4, and whose Lean status is YELLOW
  by design pending Mathlib upgrades.
\item No \textbf{empirical observation} or \textbf{open problem} is
  formalized or referenced by TS6.
\end{itemize}

TS6 does not close the TS programme: the asymptotic law for $\AWs$ remains
the programme's principal open question. TS6 locks the exact-kernel
identities at the level of typed Lean statements with documented proof
plans, producing a certified index of what the TS programme's formal
backbone looks like.

\end{document}
